\exercises

\tierunderstand

\begin{bt}
  \item Lấy tích phân~\eqref{eq:ch04-rang-buoc} lại từ đầu và chỉ ra chính xác
    chỗ nào trong phép tính buộc \(f_j\) phải tuyến tính. Rồi trả lời: nếu bỏ
    chữ \enquote{với mọi \(t\)} khỏi ràng buộc, kết luận còn đúng không?

  \item Mục~\ref{sec:ch04-dong-loi-khong-doi} đo được \(w = 4\) cho tích nhỏ
    hơn \(w = 2\) khoảng năm mươi ba bậc độ lớn, dù 4 mới là trọng số nhỏ nhất
    về lý thuyết đạt được hệ số 1. Giải thích bằng quỹ đạo của \(a_j(t)\), rồi
    dự đoán chuyện gì xảy ra ở \(w = 3\).

  \item Dẫn lại~\eqref{eq:ch04-nghiem-toi-uu} từ đầu. Rồi sửa bài toán: target
    là \(x_1 + x_2\) thay vì \(x_1\), với hai vị trí cố định. Trọng số tối ưu
    bây giờ là bao nhiêu, và phần phương sai giải thích được là bao nhiêu?

  \item Trong~\eqref{eq:ch04-jac-day-du}, ma trận \(\mat{D}\) chứa
    \(\vect{o}_{t-1}\) chứ không phải \(\vect{o}_t\). Viết ra số hạng thừa mà
    bạn sẽ thêm vào nếu chép nhầm thành \(\vect{o}_t\), và nói vì sao nó không
    thuộc về đó.

  \item Mục~\ref{sec:ch04-dao-ham} nói tín hiệu lỗi bị co ở hai đầu chứ không
    co ở giữa. Viết ra hai thừa số ấy, rồi lập luận xem điều gì xảy ra với
    chúng khi \tn{cổng ra}{output gate} bị đẩy về không: đường dài có bị ảnh
    hưởng khác đường ngắn không?

  \item Nhìn hình~\ref{fig:ch04-dong-loi}. Đường đứt đi lên là ô nhớ bỏ phép
    cắt gradient. Nếu bạn chỉ được xem đường ấy mà không xem hai đường kia, bạn
    sẽ kết luận gì về kiến trúc này, và kết luận đó sai ở đâu?

  \item Bảng tham số ở mục~\ref{sec:ch04-con-lai-gi} cho ba tỷ số: 3, 4 và 9.
    Với mỗi con số, nói rõ nó đếm kiến trúc nào và trên topology nào. Rồi trả
    lời: vì sao GRU cũng là 3 mà lại rẻ hơn LSTM 1997 về mặt tính toán?
\end{bt}

\tierapply

Trạng thái xuất phát cho cả năm bài: clone repo companion, đứng ở thư mục gốc
của nó, và checkout đúng tag của chương này.

\begin{minted}{console}
$ git clone https://github.com/Giang-Dang/rnn-to-transformer-lab.git
$ cd rnn-to-transformer-lab
$ git checkout ch04
$ conda env create -f environment.yml
$ conda activate rnn-to-transformer-lab
$ python verify.py --only ch04
\end{minted}

Dòng cuối phải in ra \code{verify: ok}. Nếu không, đừng làm tiếp: mọi con số
bạn đo sau đó sẽ không so được với sách.

\begin{bt}
  \item \code{cec\_jacobian\_truncated} trả về ma trận đơn vị và không nhận
    tham số trạng thái nào, nên \code{ch04\_flow.py} không dùng autograd ở đâu
    cả. Hệ quả: detach thêm trạng thái \tn{ô nhớ}{memory cell} trong
    \code{unroll} không đổi được một chữ số nào của bảng ấy. Hãy kiểm điều đó,
    rồi dựng phép đo mà nó \emph{có} đổi: chạy
    \code{ch04\_truncation.py} với \code{states[-1].c} bị detach thêm, và nói
    xem cosine ở \(T = 100\) đi về đâu và vì sao.

  \item \code{experiments/ch04\_truncation.py} đo cosine ở trọng số khởi tạo
    trong \([-0.1, 0.1]\). Chạy lại ở \code{SCALE = 0.02} và \code{SCALE = 0.5}
    rồi vẽ cosine theo \(T\) cho cả ba thang. Cosine tụt nhanh hơn hay chậm hơn
    khi trọng số nhỏ đi, và điều đó khớp với số đo ở
    mục~\ref{sec:ch04-dao-ham} thế nào?

  \item \code{LstmForget} trong \code{lstm.py} có \tn{cổng quên}{forget gate}
    nhưng chưa có
    experiment nào dùng nó. Viết một script dựng lại đo 3 của
    mục~\ref{sec:ch04-dao-ham} cho nó: đo
    \(\norm{\jac{\vect{c}_t}{\vect{c}_k}}\) theo khoảng cách với bias cổng quên
    đặt lần lượt ở \(-2\), \(0\) và \(+2\). Bias nào cho đường gần nhất với
    đường liền trong hình~\ref{fig:ch04-dong-loi}, và vì sao?

  \item \code{experiments/ch04\_adding.py} chạy ở \(T = 100\). Đổi
    \code{T\_LAG} thành 200 và 400, giữ nguyên mọi thứ khác, và xem ô nhớ còn
    giải được không. Ghi lại thời gian chạy: bạn sẽ chạm ngân sách 60 giây
    trước khi chạm giới hạn của kiến trúc, nên hãy nói rõ bạn đã đổi gì để đo
    tiếp.

  \item Bài 1997 dùng \(g_{\mathrm{in}}\) khoảng \([-2,2]\) và
    \(g_{\mathrm{out}}\) khoảng \([-1,1]\). Thay cả hai bằng \code{torch.tanh}
    trong \code{lstm.py}, chạy \code{python -m pytest -q tests/test\_ch04.py},
    và liệt kê test nào hỏng. Rồi chạy lại \code{ch04\_adding.py}: kết luận của
    chương có đổi không, và nếu không thì test hỏng kia đang bảo vệ cái gì?
\end{bt}

\tierextend

\begin{bt}
  \item Mục~\ref{sec:ch04-cat-gradient} lập luận rằng phép cắt gradient vứt đi
    đúng thành phần phình theo khoảng cách. Lập luận ấy dựa trên hai bảng đo hai
    đại lượng khác nhau và ghép lại bằng lời. Thiết kế một phép đo duy nhất
    kiểm trực tiếp được nó, chạy, và nói xem nó ủng hộ hay bác bỏ cách đọc của
    tôi.

  \item Bài nói bản cắt cụt không giải được bài kiểu XOR trễ, vì bài toán không
    phân rã được thành mục tiêu con dễ hơn. Dựng bài toán ấy trên bộ sinh của
    \code{ch04\_adding.py}, huấn luyện, và xem nó hỏng. Rồi tìm một cách sửa
    \emph{bài toán}, không sửa kiến trúc, làm nó phân rã được. Bạn vừa thêm giả
    định gì?

  \item Cổng quên năm 2000 đổi \(\jac{\vect{c}_t}{\vect{c}_{t-1}}\) từ
    \(\mat{I}\) thành \(\diag(\vect{f}_t)\), tức trả lại chính xác cái dạng
    tích mà chương~\ref{ch:phan-tich} chứng minh là nguy hiểm. Vậy mà nó thắng.
    Tìm điều kiện lên \(\vect{f}_t\) mà dưới đó tích không tắt, rồi xem một
    mạng đã huấn luyện xong có thoả điều kiện ấy không. Đây là một chỗ còn mở
    thật sự: chưa có phát biểu nào được chấp nhận rộng rãi về việc cổng quên
    học được cần giữ những gì để trí nhớ dài hạn sống sót.

  \item Chương này bỏ qua một câu hỏi: ô nhớ nhớ được bao nhiêu \emph{bit}.
    CEC là một số thực và float64 có 53 bit mantissa, nhưng
    \(\vect{c}_t\) còn phải chịu nhiễu từ mọi lần \tn{cổng vào}{input gate} hé
    mở. Nghĩ ra một
    cách đo dung lượng nhớ thực tế của một cell theo số bước, đo nó, rồi so với
    số cell mà bài 1997 dùng cho bài toán 2c ở quãng 1000 bước.
\end{bt}
